\section{Related Work}
\label{sec:related}

We situate our work within three streams: equivariant weight-space encoders, plain weight-space learning with augmentation, and the theoretical expressivity of permutation-equivariant networks.

\subsection{Equivariant Weight-Space Encoders}

The observation that neural network weights carry a permutation symmetry motivates encoder architectures that respect this symmetry by construction. \textbf{DWSNets} \cite{Navon2023Equivariant} introduced the first permutation-equivariant layers for homogeneous MLP weight spaces, with Navon et al.\ reporting R$^2{\approx}0.89$ for accuracy prediction on a private MNIST model zoo. \textbf{GNN-NFN} \cite{Kofinas2024Graph} extended this idea to diverse architectures by treating neural networks as computational graphs (nodes = neurons, edges = weight matrices) and applying graph neural network operations equivariant to each layer's permutation group. GNN-NFN achieves state-of-the-art property prediction across multiple architecture families. \textbf{NFN} \cite{Zhou2023Permutation} and its universal extension \cite{Zhou2024Universal} auto-construct equivariant layers for arbitrary weight spaces. \textbf{Monomial-NFN} \cite{Tran2024Monomial} further extends symmetry handling to scaling and sign-flipping, achieving 34\% parameter reduction via theory-guided construction.

However, none of these works conduct a controlled comparison with plain encoders on shared train/test splits, and none report sample efficiency curves across systematically varied training set sizes. Their performance is measured at full-data scale on custom or private splits.

\subsection{Plain Weight-Space Learning and Augmentation}

The ModelZooDataset \cite{Schurholt2022Model} provides the standardized model zoos used in our study. Building on this resource, \citet{Schurholt2021Hyper} demonstrated that self-supervised representation learning on flattened weights can predict model accuracy (R$^2{\approx}0.83$) on the ModelZooDataset MNIST zoo---establishing plain encoders as competitive baselines. \citet{Meynent2025Structure} showed that combining behavioral and structural signals outperforms structure alone.

Permutation augmentation---randomly permuting neuron orderings during training---is analogous to random crops in image learning: it achieves some but not all benefits of architectural invariance. Its role as an intermediate condition has been discussed conceptually but never tested in a controlled weight-space learning study.

These plain encoder works share a key limitation: no equivariant encoder comparison on the same shared splits.

\subsection{Expressivity Theory and Sample Complexity}

\citet{Dayan2026Expressive} proved that all permutation-equivariant weight-space networks are equivalent in expressivity. Their theorem implies equivariant and plain encoders share the same function class ceiling when data is abundant. However, expressivity equivalence does not imply sample complexity equivalence---a smaller, symmetry-consistent hypothesis space can be traversed with fewer examples even if ultimately equally expressive, analogous to why CNNs outperform plain MLPs on images at finite data.

Our work provides the empirical counterpart: we measure the practical sample complexity gap and show it closes at full data (Section~\ref{sec:results}). \citet{Herrmann2024Learning} contributed the first RNN model zoo datasets, finding functionalist representations outperform mechanistic ones---paralleling our finding that data-level augmentation can outperform structural encoding at very small zoo sizes.

\subsection{Positioning}

Our work is the first to: (1) compare equivariant and plain weight-space encoders on shared ModelZooDataset splits, (2) measure sample efficiency via systematic training size ablation, (3) include permutation augmentation as an explicit intermediate condition, and (4) report a data-regime crossover that qualifies the conventional wisdom that equivariant inductive bias helps most at low data.
